\section{Numerical evidence}\label{sec: evidence}
\subsection{Consistency checks}\label{subsec: checks}

\paragraph{(a) No foreign rate.} If $\sigma^f_r=0$ then $\zeta^f=\zeta^d=\ell^f\equiv 0$ and the first line of (\ref{eq: sensitivities}) is (\ref{eq: oneProcess}): formula (\ref{eq: mainFormula}) is the one-factor Hull-White formula (\ref{eq: malliavinHullWhite}) of \cite{guyonLabordere}.

\paragraph{(b) Constant spot volatility.} If $\sigma(t,S)\equiv\sigma$ then $g=0$, the system (\ref{eq: sensitivities}) is deterministic and can be integrated:
$$ 2\big(\ell^d_t-\ell^f_t\big)=\partial_t\int_0^t\Big[
(\sigma^d_r A^d)^2+(\sigma^f_r A^f)^2+2\rho_{Sd}\sigma\sigma^d_r A^d-2\rho_{Sf}\sigma\sigma^f_r A^f
-2\rho_{df}\sigma^d_r\sigma^f_r A^dA^f\Big](s,t)\,ds . $$
The integral is the part of the total variance of the forward $f$ that is due to the rates in the Gaussian three-factor model, so the right hand side is exactly $\sigma_{Dup}(t)^2-\sigma^2$. With zero correlations this reads
$$ \sigma^2=\sigma_{Dup}(t)^2-\frac{(\sigma^d_r)^2}{\kappa_d^2}(1-e^{-\kappa_dt})^2-\frac{(\sigma^f_r)^2}{\kappa_f^2}(1-e^{-\kappa_ft})^2 . $$
This fixes all signs and the normalisation $2/K$ of both terms in (\ref{eq: threeFactorGeneralMalliavin}).

\paragraph{(c) Monte Carlo, with skew.} Integrating (\ref{eq: mainFormula}) twice in the strike gives two identities between unconditional expectations, which can be tested without any kernel estimate: with $c^i_T=r^i_T-r^0_i(T)$,
\begin{align*}
\mathrm{(I)}\quad&\mathbb{E}^d\big[D^d_{0T}\,c^d_T\,(S_T-K)^+\big]=\mathbb{E}^d\big[D^d_{0T}\,\mathbb{I}_{S_T>K}\,S_T\,\ell^d_T\big],\\
\mathrm{(II)}\quad&\mathbb{E}^d\big[D^d_{0T}\,S_T\,c^f_T\,(S_T-K)^+\big]=\mathbb{E}^d\big[D^d_{0T}\,\mathbb{I}_{S_T>K}\,S_T^2\,\ell^f_T\big].
\end{align*}
They were tested for a fixed local volatility with a pronounced skew,
$$\sigma(S)=0.12\,\big(1-0.6\tanh(2\ln S)\big),\qquad S_0=1,\qquad T=5,$$
with $\kappa_d=5\%$, $\sigma^d_r=1.2\%$, $\kappa_f=10\%$, $\sigma^f_r=1.5\%$, correlations
$\rho_{Sd}=0.35$, $\rho_{Sf}=-0.30$, $\rho_{df}=0.40$, flat curves at $3\%$ and
$1\%$, $2\cdot10^6$ antithetic paths and 200 Euler steps per year. In the table, $z$ is the difference between the two sides in standard errors. The column ``naive'' uses the system (\ref{eq: sensitivities}) without the cross terms and with $g(dW^S-\sigma dt)$ in $\ell^f$, i.e. without the corrections due to the measure change.

\begin{center}
\begin{tabular}{c|rrr|rrr|rr}
\toprule
 & \multicolumn{3}{c|}{(I) domestic, $\times10^{-3}$} & \multicolumn{5}{c}{(II) foreign, $\times10^{-3}$}\\
$K$ & direct & (\ref{eq: sensitivities}) & $z$ & direct & (\ref{eq: sensitivities}) & $z$ & naive & $z$\\
\midrule
0.70 & $2.4246$ & $2.4219$ & $+0.42$ & $-3.0956$ & $-3.0939$ & $-0.15$ & $-3.1798$ & $+6.96$\\
0.85 & $2.1273$ & $2.1244$ & $+0.52$ & $-2.8555$ & $-2.8537$ & $-0.18$ & $-2.9303$ & $+7.19$\\
1.00 & $1.7123$ & $1.7086$ & $+0.80$ & $-2.4536$ & $-2.4512$ & $-0.28$ & $-2.5123$ & $+6.98$\\
1.15 & $1.2363$ & $1.2361$ & $+0.04$ & $-1.9139$ & $-1.9138$ & $0.00$ & $-1.9562$ & $+6.66$\\
1.30 & $0.7868$ & $0.7876$ & $-0.35$ & $-1.3241$ & $-1.3245$ & $+0.08$ & $-1.3497$ & $+5.48$\\
1.50 & $0.3491$ & $0.3497$ & $-0.46$ & $-0.6578$ & $-0.6584$ & $+0.21$ & $-0.6684$ & $+3.65$\\
\bottomrule
\end{tabular}
\end{center}

\noindent Both identities hold within one standard error, while the naive foreign term is off by about $2.7\%$. For these parameters the cross term in the domestic pair is small (about $0.2\%$, not resolved by the test). In a stressed configuration ($T=10$, both rate volatilities doubled, skew coefficient $0.8$, 100 steps per year) both identities again hold within one standard error, the naive domestic term is off by about $1\%$ ($z$ between $3$ and $5$) and the naive foreign term by about $3\%$ ($z$ between $4$ and $8$).

The simulation scheme of Section \ref{subsec: discretisation3F} was checked separately: the closed forms of all step coefficients, including $K^{df}_j$, agree with numerical quadrature to $10^{-10}$; under $\mathbb{Q}^{d^T}$ the simulated particles reproduce $\mathbb{E}[f_t]=f_0$, $P^d_{0T}\mathbb{E}[(P^d_{tT})^{-1}]=P^d_{0t}$ and $P^d_{0T}\mathbb{E}[(P^d_{tT})^{-1}S_t]=S_0P^f_{0t}$ within Monte Carlo error; and (I)--(II) hold at an intermediate date ($t=4$, $T=6$) with $\mathbb{Q}^{d^T}$ particles and weights $(P^d_{tT})^{-1}$. Pooling five independent runs of $10^6$ paths at that date (two under $\mathbb{Q}^{d^T}$, three under $\mathbb{Q}^d$), the two sides differ at $K=0.7$ by $0.13\%$ for (I) and $0.09\%$ for (II), that is $0.8$ and $0.4$ standard errors; single runs deviate by up to $2.9$ standard errors at one strike, and the right hand sides computed under the two measures agree to about $0.1\%$.
The matrix form (\ref{eq: matrixSystem}), simulated directly, agrees with (\ref{eq: sensitivities}) path by path up to the Euler error. Scripts: \texttt{check\_sensitivities.py}, \texttt{check\_forward\_measure.py}, \texttt{check\_matrix\_form.py}.

\subsection{Calibration to a market surface}\label{subsec: cme}
The Monte Carlo evidence of Section \ref{subsec: checks} concerns the identities behind (\ref{eq: mainFormula}) for a \emph{given} smooth local volatility, and the simulation scheme. The full particle calibration has been run once, on CME EUR/USD option settlements of 12 November 2026: expiries from December 2026 to December 2027, the surface being extended beyond that at constant implied-volatility smile to five years, which is synthetic. The calibration is run under the domestic risk-neutral measure with discount weights, 200,000 particles and 100 time steps per year, and two sets of rate parameters which are illustrative and not calibrated to interest rate options. The options are then repriced with an independent simulation. Over the quoted expiries the repricing error attributable to the stochastic rates is at most 0.005 volatility points with the correction (\ref{eq: mainFormula}), against 0.19 and 0.47 for the two parameter sets with the Dupire local volatility alone. With the more volatile set, the calibrated variance reaches its floor within one standard deviation of the money from about three years, on the synthetic part of the surface, and the five-year at-the-money error is 0.54. The report and the scripts are in the repository \texttt{github.com/ceikit/sir-lsv-malliavin-particle}, file \texttt{calibration/RESULTS.md}.

\subsection{Synthetic comparison of three methods}\label{subsec: synthetic}
A second test is synthetic and uses options on the spot. A non-flat local volatility, whose level and skew vary in time, generates the target Dupire variance on a two-year horizon. The volatility is then recovered from that target by three methods: a three-dimensional forward equation for the discounted density with the conventional tail integrals, the five-field system of Section \ref{subsec: momentClosure}, and the particle method with the sensitivities (\ref{eq: sensitivities}). On their finest settings the three methods recover the generating volatility with a root-mean-square error between 0.002 and 0.003 volatility points, and the calibrated models reprice options against the generating model within 0.004. With 20,000 particles the particle method has an error of 0.004 after about three seconds; the two PDE solvers take 69 and 189 seconds on their finest grids. This is one set of parameters and one implementation of each method, and at this accuracy the error of the reference target is not negligible, so the methods cannot be ranked. Details are in \texttt{benchmarks/RESULTS.md}.

\subsection{What remains untested}
Several things remain untested. The CME options are options on futures and were treated as options on the forward to their expiry. The bias due to the time step and to the kernel bandwidth has not been quantified systematically. No market data beyond December 2027 was used, and no comparison has been made with other particle methods such as \cite{cozma,ogetbilHientzsch}.
